\usepackage[dvipsnames]{xcolor} 
\usepackage{soul}

%% Comentários do professor
\usepackage[textwidth=18mm,colorinlistoftodos,prependcaption,textsize=tiny]{todonotes}

\setlength{\marginparwidth}{19mm}

\usepackage{xargs} 
\newcommandx{\incerto}[2][1=]{\todo[linecolor=orange,backgroundcolor=orange!25,bordercolor=orange,#1]{#2}}
\newcommandx{\alterar}[2][1=]{\todo[linecolor=blue,backgroundcolor=blue!25,bordercolor=blue,#1]{#2}}
\newcommandx{\info}[2][1=]{\todo[linecolor=green,backgroundcolor=green!25,bordercolor=green,#1]{#2}}
\newcommandx{\melhorar}[2][1=]{\todo[linecolor=yellow,backgroundcolor=yellow!25,bordercolor=yellow,#1]{#2}}

%marcação das correções
\newcommand{\revisadoprofessor}[1]{%
    \todo[linecolor=RoyalBlue,backgroundcolor=RoyalBlue!25,bordercolor=RoyalBlue,inline]{ %
        \textbf{Revisado até aqui pela professora Angelita.} \vspace{0.3em} \newline
        \textbf{Orientações:} \newline
        - Faça os ajustes necessários. \newline
        - Para ajustes menores, apenas realize o ajuste e retire as marcações da professora. \newline
        - Para ajustes maiores, use a tag \texttt{\textbackslash corrigido\{texto corrigido\}}, abrindo parênteses no início da correção e fechando no final.  Você pode incluir múltiplos parágrafos dentro da tag.\newline
        #1
    }%
}

\long\def\corrigido#1{{\color{RoyalBlue}#1}}
% \corrigido{
% Este é um texto longo que a professora pediu para alterar. 
% Ele pode ter vários parágrafos.

% Como este aqui, tudo continuará em azul para facilitar as próximas correções.
% }

% Para riscar o texto
\newcommand{\remover}[1]{%
  \setstcolor{red}\st{#1}
}

% Comentario no texto
\newcommand{\add}[1]{%
  \textcolor{RoyalBlue}{\textbf{#1}}
}

% Complementar frase
\newcommand{\complemento}[1]{%
  \textcolor{red}{???}\todo[linecolor=red,backgroundcolor=red!25,bordercolor=red]{Falta complemento: {#1}}
}

%incluir citação
\newcommand{\referencia}{%
  \textcolor{red}{\(Ref!\)}\todo[linecolor=red,backgroundcolor=red!25,bordercolor=red]{Incluir referencia aqui}
}

\newcommand{\virgula}{%
  \textcolor{red}{\textbf{,}}
}

\newcommand{\substituir}[2]{%
    \setstcolor{red}\st{#1}
    \textcolor{red}{#2}
}

%quebrar frase/parágrafo
\newcommand{\breakfr}{%
  \textcolor{red}{\textbf{|}}\todo[linecolor=red,backgroundcolor=red!25,bordercolor=red]{QUEBRA DE FRASE}
}
\newcommand{\breakpr}{%
  \textcolor{red}{$\hookleftarrow$} \todo[linecolor=red,backgroundcolor=red!25,bordercolor=red]{QUEBRA DE FRASE}
}

\newcommand{\corrigir}[2]{%
    \hl{#1}%
    \todo[linecolor=orange,backgroundcolor=orange!25,bordercolor=orange]{#2}%
}

%comando grifar
\usepackage{tikz}
\usetikzlibrary{shapes.geometric}
\newcommand{\grifar}[1]{%
  \begin{tikzpicture}[baseline=(char.base)]
    \node[draw=red, ellipse, inner sep=2pt, thick] (char) {#1};
  \end{tikzpicture}%
}
